\documentclass[10pt,a4paper]{amsart}
\usepackage[margin=19mm]{geometry}
\usepackage{amsmath,amssymb,mathtools}
\usepackage[scr=rsfs,cal=euler]{mathalfa}
\usepackage{xfrac}
\usepackage[unicode,hidelinks,bookmarks=false]{hyperref}
\newcommand{\ie}{i.e.\ }
\newcommand{\cf}{cf.\ }
\newcommand{\resp}{respectively\ }
\newcommand{\Subsection}{\subsection}
\providecommand{\nobreakdash}{\nobreak\hbox{-}}
\setlength{\emergencystretch}{2em}
\multlinegap0pt
\usepackage{bm}
\usepackage{bbm}
\usepackage{dsfont}
\usepackage{xspace}
\usepackage{cancel}
\usepackage{mathtools}
\usepackage{amsthm}
\makeatletter
\@ifundefined{lemma}{\newtheorem{lemma}{Lemma}}{}
\@ifundefined{theorem}{\newtheorem{theorem}{Theorem}}{}
\makeatother
\DeclareMathOperator{\rad}{rad}
\newcommand{\Li}[2]{\left(\frac{#1}{#2}\right)}
\newcommand{\N}{\mathbb{N}} 
\newcommand{\OO}{\mathcal{O}}
\newcommand{\Z}{\mathbb{Z}}
\newcommand{\pp}{\mathfrak{p} }
\title[On Dold condition and fail factor of linear recurrent sequen]{On Dold condition and fail factor of linear recurrent sequences}
\author{Mateusz Rajs}
\date{Source: arXiv:2509.09847v1 (CC BY 4.0)}
\begin{document}
\maketitle

\noindent\emph{Source and licence.} On Dold condition and fail factor of linear recurrent sequences. Version arXiv:2509.09847v1 (CC BY 4.0), distributed under the Creative Commons Attribution 4.0 International licence, as recorded in the arXiv licence metadata for this exact version and shown on the abstract page https://arxiv.org/abs/2509.09847v1. Full rights notice: licenses/arxiv-2509.09847-CC-BY-4.0.txt.\par\smallskip

\noindent\textit{Editorial scope.} Counted unit: the Theorem whose source label is \texttt{TheoremNotSquare} in the article (source0.tex lines 400--408, source label \texttt{TheoremNotSquare}), delivered together with the article's own complete original proof of it (source0.tex lines 411--523, one \texttt{proof} environment of 113 lines). No mathematics is written, repaired or re-derived: statement and proof are the article's own text line for line. Required context retained uncounted: lem\_allInOne (lines 323--328). Every definition, display and label of those lines is retained unchanged. Omitted from this package and not redistributed: 1--322, 329--399, 409--410, 524--2029. Nothing inside a retained excerpt is omitted or altered: every retained body is delivered exactly as the article's lines are written, so no omission receipt of any kind accompanies this package. Citations: the retained excerpts carry no citation command at all, so no bibliography entry of the article is transcribed and no reference is redistributed. Reference adaptation: none was needed and none was made; every reference inside the delivered unit resolves inside it, so no unresolved reference or citation is produced. Printed slips kept literal: no printed word, symbol, hypothesis or number of the counted statement or of its proof was altered. Layout and class: the article's own class is \texttt{elsarticle} with packages including amssymb, amsmath, amsthm, amsfonts, inputenc, xcolor, makecell; the wrapper is the lane's pinned \texttt{amsart} editorial wrapper at 10pt on A4, which restates the theorem-like environments the retained text needs and adds the article's own macro definitions that the retained text needs (\texttt{\textbackslash Li}, \texttt{\textbackslash N}, \texttt{\textbackslash OO}, \texttt{\textbackslash Z}, \texttt{\textbackslash pp}, \texttt{\textbackslash rad}), with extra wrapper packages (bm, bbm, dsfont, xspace, cancel, mathtools). The delivered numbering is the unit's own. Assets: no retained span contains a figure environment, a graphics-inclusion command, a PDF-page inclusion, a table or a TikZ picture; no image file is copied into this package, the package declares no asset dependency and no image file is redistributed with it. Licence: version arXiv:2509.09847v1 of this article is distributed under the Creative Commons Attribution 4.0 International licence, as recorded in the arXiv licence metadata for this exact version and shown on the abstract page https://arxiv.org/abs/2509.09847v1. A full rights notice is delivered at licenses/arxiv-2509.09847-CC-BY-4.0.txt. Characters: every retained excerpt is pure ASCII, so the delivered engine, XeLaTeX with its default Latin Modern text font, reports no missing character.
\begin{lemma}\label{lem_allInOne}
    Let $K$ be a number field and $x,y\in\OO_K$. Let $\pp$ be a prime ideal of $\OO_K$ lying over a prime number $p$. Then 
    % satisfying $x\equiv y\pmod{\pp^d}$ for some $d>0$ and prime ideal $\pp$, then $x^{p^k}\equiv y^{p^k} \pmod{\pp^{d+k}}$ for every $k\geqslant0$ f
    $$ x \equiv y \pmod{\pp^{de}} \quad\Rightarrow\quad x^{\pp^k} \equiv y^{p^k} \pmod{\pp^{e(d+k)}}$$
    for $k, d\in \N_+$.
\end{lemma}

\begin{theorem}\label{TheoremNotSquare}
    The following conditions are equivalent:
    \begin{enumerate}
        \itemsep-0.2em
        \item the sequence ${\mathbf U}$ almost satisfies the Dold condition, 
        \item $l_1=l_2$, 
        \item the sequence $2r_2\rad(\Delta_\mathbf{U}){\mathbf U}$ satisfies the Dold condition. 
    \end{enumerate}
\end{theorem}

\begin{proof}

    $(1\Rightarrow 2)$

    If the sequence ${\mathbf U}$ is almost realizable, then $p|U_p-U_1$ for all but finitely many prime numbers $p$. We take a prime number $p$ such that $p|U_p-U_1$, $\left(\frac{\Delta}{p}\right)=-1$ and $p$ does not divide denominators of $l_1$ and $l_2$. There are infinitely many such primes. Notice that:
    \begin{equation}
      \begin{split}
            0 &\equiv
            U_p - U_1 \equiv
            (l_1\alpha^p + l_2\beta^p) - (l_1\alpha + l_2\beta) \\ 
            &\equiv 
            (l_1\beta + l_2\alpha) - (l_1\alpha + l_2\beta) \equiv
            (l_1-l_2)(\beta - \alpha) \pmod p
      \end{split}
      \end{equation}
    Therefore, $\alpha \equiv \beta \pmod p$ or $l_1 \equiv l_2 \pmod p$. This property holds for infinitely many prime numbers $p$. If $\alpha \equiv \beta \pmod p$ for infinitely many numbers $p$, then  $\alpha = \beta$, but then the quadratic equation has only one solution, hence $\Delta_{\mathbf U} = 0$, which contradicts the assumption that $\Delta_{\mathbf U}$ is not a square of an integer. Therefore, $l_1 \equiv l_2 \pmod p$ must be true for infinitely many prime numbers $p$, which implies that $l_1=l_2$.

    \vskip 0.5cm
    \noindent $(2\Rightarrow3)$

    The coefficients $l_1 = l_2$ may not be integers. To be able to analyze them modulo prime $p$ we must first make sure that $p$ does not divide their denominators. Fortunately, one can notice that $2r_2l_1$ is an integer because
    \begin{align}
        2r_2l_1 = r_2(l_1+l_2) = r_2U_0 = U_2 - r_1U_1 \in \Z ,
    \end{align}
    where the second equality comes from the exact formula and the third one from the recurrence relation.

    We will show that the sequence $\rad(\Delta_\mathbf{U})\mathbf{V}$ satisfies the Dold condition where ${V}_n := \alpha^n+\beta^n$. This will imply that $2r_2\rad(\Delta_\mathbf{U})\mathbf{U}$ also satisfies this condition as this sequence is an integer multiple of $\rad(\Delta_\mathbf{U})\mathbf{V}$. 

    To show that $\rad(\Delta_\mathbf{U})\mathbf{V}$ satisfies the condition we will prove that for every prime $p$ we have
    \begin{align}\label{the_not_sq_e1}
        \begin{split}
        \rad(\Delta_\mathbf{U})\left(\alpha^{sp^k}+\beta^{sp^k}\right) \equiv 
        \rad(\Delta_\mathbf{U})\left(\alpha^{sp^{k-1}}+\beta^{sp^k}\right) 
        \pmod {p^k} 
        \\ \text{for all $k, s\in \N_+$}
        \end{split}
    \end{align}

    
For any prime $p$ the expression $\Delta_\mathbf{U}^\frac{p-1}{2} \pmod p$ (and consequently $\Li{\Delta_\mathbf{U}}{p}$) can only take on one of three values. 
\begin{enumerate}
    \item $\Delta_\mathbf{U}^\frac{p-1}{2} \equiv -1 \pmod p$ and $\Delta_\mathbf{U}$ is not a quadratic residue. We already considered this case and saw that
    % \begin{align}
    %     \sqrt{\Delta_{\mathbf U}}^p \equiv
    %     \Delta_{\mathbf U}^{\frac{p-1}{2}} \cdot \sqrt{\Delta_{\mathbf U}}
    %     \pmod p.
    % \end{align}
    % Thus, for any integers $a,b$ we have
    % \begin{align}
    %      \left(a+b\sqrt{\Delta_{\mathbf U}}\right)^p \equiv
    %      a^p + b^p\sqrt{\Delta_{\mathbf U}}^p \equiv
    %      a - b\sqrt{\Delta_{\mathbf U}}
    %      \pmod{p}
    % \end{align}
    % Consequentially, the transformation $x \mapsto x^p$ permutates the roots modulo $p$, that is
    $\alpha^p\equiv\beta\pmod p$ and $\beta^p \equiv \alpha\pmod p$. By Lemma we have \ref{lem_allInOne}  $\alpha^{p^k}\equiv\beta^{p^{k-1}}\pmod {p^k}$ and $\beta^{p^k}\equiv\alpha^{p^{k-1}}\pmod {p^k}$. Taking the last two equivalences to the power of $s$ and adding them together we get (\ref{the_not_sq_e1}).

    
    \item $\Delta_\mathbf{U}^\frac{p-1}{2} \equiv 1 \pmod p$ and $\Delta_\mathbf{U}$ is a quadratic residue. This case follows similarly. We have
     \begin{align}
         \left(A+B\sqrt{\Delta_{\mathbf U}}\right)^p \equiv
         A^p + B^p\cdot \Delta_{\mathbf U}^{\frac{p-1}{2}} \cdot \sqrt{\Delta_{\mathbf U}} \equiv
         A + B\sqrt{\Delta_{\mathbf U}}
         \pmod{p}.
    \end{align}
    Therefore $\alpha^p\equiv\alpha\pmod p$ and $\beta^p \equiv \beta\pmod p$.  By Lemma \ref{lem_allInOne}  $\alpha^{p^k}\equiv\alpha^{p^{k-1}}\pmod {p^k}$ and $\beta^{p^k}\equiv\beta^{p^{k-1}}\pmod {p^k}$. Again, taking both to the power of $s$ and adding them we get $(\ref{the_not_sq_e1})$.

    \item $\Delta_\mathbf{U}^\frac{p-1}{2} \equiv 0 \pmod p$ and $p|\Delta_\mathbf{U}$. We have
     \begin{align}
         \left(A+B\sqrt{\Delta_{\mathbf U}}\right)^p \equiv
         A^p + B^p\sqrt{\Delta_{\mathbf U}}^p \equiv
         A
         \pmod{p}.
    \end{align}
    Hence there must exist $A\in\Z$ such that $\alpha^p\equiv\beta^p\equiv A\pmod p$. Again by Lemma \ref{lem_allInOne} we have $\alpha^{p^k}\equiv\beta^{p^k}\equiv A^{p^{k-1}}\pmod {p^k}$. We can see that
    \begin{align}
        \alpha^{p^ks} \equiv  
        \left(   \alpha^{p^k} \right)^s  \equiv 
         \left( A^{p^{k-1}}\right)^s  \equiv
         \left( A^{p^{k-2}}\right)^s  \equiv
         \\\equiv
         \left(\alpha^{p^{k-1}}\right)^s  \equiv
          \alpha^{p^{k-1}s} \pmod {p^{k-1}}
    \end{align}
    for any $s\in\N_+$. The second equivalence can be proven by applying Lemma \ref{lem_allInOne} to $A^p\equiv A\pmod p$ (which is true by the Fermat's little theorem). The similar argument works to show that  $\beta^{p^k} \equiv   \beta^{p^{k-1}s} \pmod {p^{k-1}}$. Adding these results and multiplying them by $p$ we get  $ p\left(\alpha^{sp^k}+\beta^{sp^k}\right) \equiv  p\left(\alpha^{sp^{k-1}}+\beta^{sp^k}\right) \pmod {p^k}$. Because in this case $p|\rad(\Delta_\mathbf{U})$ this implies (\ref{the_not_sq_e1}). 

\end{enumerate}

% If a prime $p$ falls into the first or second case then clearly  $\alpha^{p^k}+\beta^{p^k}\equiv\alpha^{p^{k-1}}+\beta^{p^{k-1}}\pmod {p^k}$ for $k\in\N_+$.

% The third case requires more work.
   


    
    
    % by the recurrence formula $U_2 - r_1U_1= r_2U_0$ and by the exact formula $r_2U_0 = r_2(l_1+l_2) = 2r_2l_1$ so $2r_2l_1$ is an integer.

    
    % Having shown that $\alpha^p\equiv \beta \pmod p$, using Lemma \ref{lem_allInOne} we get $\alpha^{p^k} \equiv \beta^{p^{k-1}} \pmod{p^k}$ and analogously $\beta^{p^k} \equiv \alpha^{p^{k-1}} \pmod{p^k}$. 
    % By the recurrence formula $U_2 - r_1U_1= r_2U_0$ and by the exact formula $r_2U_0 = r_2(l_1+l_2) = 2r_2l_1$ so $2r_2l_1$ is an integer.
    
    % We know that $U_n = l_1(\alpha^n + \beta^n)$ since $l_1=l_2$. It is easy to notice that
    % \[
    %     2r_2U_{sp^k}  \equiv
    %     (2r_2l_1)(\alpha^{sp^k} + \beta^{sp^k}) \equiv
    %     (2r_2l_1)(\beta^{sp^{k-1}} + \alpha^{sp^{k-1}}) \equiv
    %     2r_2U_{sp^{k-1}} \pmod{p^k}.
    % \]


    \noindent $(3\Rightarrow1)$ Trivial.
\end{proof}

\end{document}
